%% Title Page (Unblinded)
%% Marketing Letters -- Replication Corner
%% Manuscript: "Does AI Reduce the Reputation Premium in Knowledge Diffusion?
%%              A Conceptual Replication Using Chinese Journal Articles"

\documentclass[12pt,a4paper]{article}

\usepackage[top=2.54cm,bottom=2.54cm,left=2.54cm,right=2.54cm]{geometry}
\usepackage{setspace}
\usepackage{amsmath}
\usepackage{times}
\usepackage[utf8]{inputenc}
\usepackage[T1]{fontenc}
\usepackage{microtype}
\usepackage{parskip}

\singlespacing
\setlength{\parskip}{0.8em}

\begin{document}

% ---- Title ------------------------------------------------------------------
\begin{center}
{\large\textbf{Does AI Reduce the Reputation Premium\\[0.3em]
in Knowledge Diffusion?\\[0.3em]
A Conceptual Replication Using Chinese Journal Articles}}

\medskip
{\normalsize\textit{Replication Corner / Marketing Letters}}
\end{center}

\bigskip

% ---- Authors and Affiliations -----------------------------------------------
\noindent
[Author names and affiliations --- blinded for review]

\bigskip

% ---- Corresponding author ---------------------------------------------------
\noindent
\textbf{Corresponding author:} [Blinded for review]

\bigskip\hrule\bigskip

% ---- Abstract ---------------------------------------------------------------
\noindent\textbf{Abstract}

\medskip\noindent
Jo and Li (2026) report that AI-style writing attenuates the
institutional reputation premium in SSRN marketing working papers.
We conceptually replicate this finding using 843 peer-reviewed Chinese
business and management journal articles indexed on CNKI (2023--2026).
AI-like writing style is proxied by a dictionary score counting 23
pre-specified structural markers in abstract text; institutional prestige
is coded via a four-tier Chinese university designation scheme.
The AI-style proxy predicts greater download volume
($\hat\beta = 0.22$, $p < .001$), partially replicating Jo and Li (2026).
However, institutional prestige is not a significant predictor of
diffusion in any specification (all $p > .10$), consistent with
editorial certification reducing the marginal role of institutional
affiliation in peer-reviewed venues.
The Prestige\,$\times$\,AI-style interaction is near zero and
non-significant across all four specifications.
These results identify a boundary condition: AI can attenuate
reputation-based diffusion only where a prestige premium is operative.

\bigskip

% ---- Keywords ---------------------------------------------------------------
\noindent\textbf{Keywords:} AI writing style; institutional prestige;
knowledge diffusion; reputation premium; replication; CNKI; China

\bigskip\hrule\bigskip

% ---- Compliance with Ethical Standards --------------------------------------
\noindent\textbf{Compliance with Ethical Standards}

\medskip
\noindent\textbf{Conflict of interest:} The authors declare that they have
no conflict of interest.

\medskip
\noindent\textbf{Funding:} Not applicable.

\medskip
\noindent\textbf{Ethical approval:} Not applicable.
This study uses publicly available secondary data (CNKI article records)
and involves no collection of personal data, no human subjects research,
and no interventions requiring ethical review.

\medskip
\noindent\textbf{Informed consent:} Not applicable.

\end{document}
